\documentclass[12pt]{article}
\usepackage[UTF8]{inputenc}
\usepackage{thaisupp}
\usepackage{geometry}
\geometry{a4paper, margin=2.5cm}
\begin{document}
\begin{center}
{\LARGE\textbf{ฟิสิกส์ ม.ปลาย — Part 1}}\\[0.3em]
{\Large\textbf{Lesson 9: คลื่น}}\\[0.5em]
\rule{0.8\linewidth}{0.5pt}
\end{center}
\section*{คำถาม In-Class}
\begin{enumerate}
\item \textbf{คลื่นในน้ำตื้นเป็นคลื่นชนิดใด}
\begin{enumerate}
\item ก) คลื่นตามยาว
\item ข) คลื่นตามขวาง
\item ค) คลื่นผสม
\item ง) คลื่นแม่เหล็กไฟฟ้า
\end{enumerate}
\item \textbf{เมื่อคลื่นเคลื่อนที่จากตัวกลางหนึ่งไปยังอีกตัวกลางหนึ่ง �ข้อใดไม่เปลี่ยนแปลง}
\begin{enumerate}
\item ก) ความยาวคลื่น
\item ข) ความถี่
\item ค) อัตราเร็ว
\item ง) ทิศการเคลื่อนที่
\end{enumerate}
\item \textbf{คลื่นเสียงเป็นคลื่นชนิดใด}
\begin{enumerate}
\item ก) คลื่นตามยาว
\item ข) คลื่นตามขวาง
\item ค) คลื่นผิวน้ำ
\item ง) คลื่นแม่เหล็กไฟฟ้า
\end{enumerate}
\item \textbf{ปรากฏการณ์ที่คลื่นเบนออกจากแนวตรงเมื่อผ่านช่องแคบเรียกว่าอะไร}
\begin{enumerate}
\item ก) การสะท้อน
\item ข) การหักเห
\item ค) การเลี้ยวเบน
\item ง) การแทรกสอด
\end{enumerate}
\item \textbf{เมื่อคลื่นตกกระทบผิวสะท้อน มุมตกกระทบเท่ากับมุมสะท้อน นี่คือกฎของฟิสิกส์ข้อใด}
\begin{enumerate}
\item ก) กฎการเลี้ยวเบน
\item ข) กฎการสะท้อน
\item ค) กฎสเนลล์
\item ง) หลักฮอยเกนส์
\end{enumerate}
\item \textbf{คลื่นน้ำเคลื่อนที่ด้วยอัตราเร็ว 4 m/s มีความยาวคลื่น 0.5 m ความถี่ของคลื่นเป็นเท่าใด}
\begin{enumerate}
\item ก) 2 Hz
\item ข) 4 Hz
\item ค) 8 Hz
\item ง) 0.125 Hz
\end{enumerate}
\item \textbf{คลื่นเสียงเคลื่อนที่ในอากาศด้วยอัตราเร็ว 340 m/s ที่ความถี่ 170 Hz จะมีความยาวคลื่นเท่าใด}
\begin{enumerate}
\item ก) 0.5 m
\item ข) 1 m
\item ค) 2 m
\item ง) 4 m
\end{enumerate}
\item \textbf{เมื่อคลื่นเคลื่อนที่จากน้ำลึกเข้าสู่น้ำตื้น ข้อใดเกิดขึ้น}
\begin{enumerate}
\item ก) ความถี่เพิ่มขึ้น
\item ข) ความยาวคลื่นเพิ่มขึ้น
\item ค) อัตราเร็วลดลง
\item ง) อัตราเร็วเพิ่มขึ้น
\end{enumerate}
\item \textbf{ช่องแคบกว้าง 2 m มีคลื่นความยาวคลื่น 0.5 m ผ่านเข้ามา จะเกิดการเลี้ยวเบนอย่างชัดเจนหรือไม่}
\begin{enumerate}
\item ก) เกิดอย่างชัดเจน เพราะ λ > ช่องแคบ
\item ข) เกิดอย่างชัดเจน เพราะ λ < ช่องแคบ
\item ค) ไม่เกิดการเลี้ยวเบน เพราะ λ > ช่องแคบ
\item ง) ไม่เกิดการเลี้ยวเบน เพราะ λ < ช่องแคบ
\end{enumerate}
\item \textbf{คลื่นสองขบวนมาจากแหล่งกำเนิดต่างกัน มีความถี่เท่ากันแต่เฟสต่างกัน พบบริเวณที่คลื่นแทรกสอดและกระทำกันแบบเสริม เรียกบริเวณนั้นว่าอะไร}
\begin{enumerate}
\item ก) บริเวณปฏิบัพ
\item ข) บริเวณบัพ
\item ค) บริเวณคงที่
\item ง) บริเวณเลี้ยวเบน
\end{enumerate}
\item \textbf{สิ่งที่คลื่นน้ำพาเคลื่อนที่ผ่านไปคือข้อใด}
\begin{enumerate}
\item ก) ตัวกลาง (น้ำ)
\item ข) พลังงาน
\item ค) ทั้งตัวกลางและพลังงาน
\item ง) ไม่พา什么都没有
\end{enumerate}
\item \textbf{เมื่อแสงเคลื่อนที่จากอากาศเข้าสู่แก้ว ข้อใดเกิดขึ้น}
\begin{enumerate}
\item ก) ความถี่เพิ่มขึ้น
\item ข) ความถี่ลดลง
\item ค) ความยาวคลื่นเพิ่มขึ้น
\item ง) ความยาวคลื่นลดลง
\end{enumerate}
\item \textbf{ปรากฏการณ์ที่คลื่นสองขบวนมารวมกันแล้วเกิดแอมพลิจูดรวมมากขึ้นหรือน้อยลง เรียกว่าอะไร}
\begin{enumerate}
\item ก) การเลี้ยวเบน
\item ข) การสะท้อน
\item ค) การแทรกสอด
\item ง) การหักเห
\end{enumerate}
\item \textbf{คลื่นเสียงเดินทางในของแข็งเร็วกว่าในของเหลว และในของเหลวเร็วกว่าในแก๊ส เพราะเหตุใด}
\begin{enumerate}
\item ก) ความหนาแน่นต่างกัน
\item ข) ระยะห่างระหว่างโมเลกุลต่างกัน
\item ค) อุณหภูมิต่างกัน
\item ง) ความดันต่างกัน
\end{enumerate}
\item \textbf{คลื่นความถี่ 50 Hz มีความยาวคลื่น 4 m คลื่นเคลื่อนที่ด้วยอัตราเร็วเท่าใด}
\begin{enumerate}
\item ก) 12.5 m/s
\item ข) 46 m/s
\item ค) 54 m/s
\item ง) 200 m/s
\end{enumerate}
\item \textbf{สมบัติของคลื่น 4 ข้อที่แสดงว่าคลื่นเป็นปรากฏการณ์ทางฟิสิกส์ประกอบด้วยข้อใด}
\begin{enumerate}
\item ก) สะท้อน หักเห แทรกสอด เลี้ยวเบน
\item ข) สะท้อน หักเห ดูดกลืน เลี้ยวเบน
\item ค) หักเห แทรกสอด เลี้ยวเบน กระเจิง
\item ง) สะท้อน แทรกสอด เลี้ยวเบน ดูดกลืน
\end{enumerate}
\item \textbf{ในการทดลองช่องคู่ของยัง ระยะระหว่างช่อง 0.5 mm ระยะจากช่องถึงฉาก 2 m ความยาวคลื่น 600 nm ระยะระหว่างปฏิบัพที่อยู่ติดกันเป็นเท่าใด}
\begin{enumerate}
\item ก) 0.6 mm
\item ข) 1.2 mm
\item ค) 2.4 mm
\item ง) 4.8 mm
\end{enumerate}
\item \textbf{คลื่นเสียงสามารถเดินทางผ่านสุญญากาศได้หรือไม่}
\begin{enumerate}
\item ก) ได้ เพราะเป็นคลื่นแม่เหล็กไฟฟ้า
\item ข) ได้ เพราะเป็นคลื่นกล
\item ค) ไม่ได้ เพราะต้องอาศัยตัวกลาง
\item ง) ได้ ขึ้นอยู่กับความถี่
\end{enumerate}
\item \textbf{ข้อใดต่อไปนี้ไม่ใช่คลื่นกล}
\begin{enumerate}
\item ก) คลื่นเสียง
\item ข) คลื่นในสปริง
\item ค) คลื่นแม่เหล็กไฟฟ้า
\item ง) คลื่นผิวน้ำ
\end{enumerate}
\item \textbf{คลื่นตามยาวและคลื่นตามขวางแตกต่างกันอย่างไร}
\begin{enumerate}
\item ก) ต่างกันที่ทิศสั่นของอนุภาคเทียบกับทิศการเคลื่อนที่ของคลื่น
\item ข) ต่างกันที่อัตราเร็วของคลื่น
\item ค) ต่างกันที่ความถี่ของคลื่น
\item ง) ต่างกันที่ความยาวคลื่น
\end{enumerate}
\item \textbf{คลื่นน้ำเคลื่อนที่จากน้ำตื้นเข้าสู่น้ำลึก โดยมีมุมตกกระทบ 30° และมุมหักเห 45° อัตราเร็วในน้ำตื้นเป็นเท่าใด (อัตราเร็วในน้ำลึก = 30 m/s)}
\begin{enumerate}
\item ก) 14.1 m/s
\item ข) 21.2 m/s
\item ค) 42.4 m/s
\item ง) 60 m/s
\end{enumerate}
\item \textbf{เมื่อคลื่นผิวน้ำกระทบสิ่งกีดขวางที่มีช่องเปิด จะเกิดปรากฏการณ์ใด}
\begin{enumerate}
\item ก) การสะท้อน
\item ข) การหักเห
\item ค) การแทรกสอด
\item ง) การเลี้ยวเบน
\end{enumerate}
\item \textbf{ในการทดลองแทรกสอดของคลื่นน้ำ จากช่องคู่ ระยะระหว่างปฏิบัพที่ 0 ถึงปฏิบัพที่ 5 บนฉายวัดได้ 20 cm ระยะจากช่องถึงฉาย 50 cm ความยาวคลื่นเป็นเท่าใด}
\begin{enumerate}
\item ก) 2 cm
\item ข) 4 cm
\item ค) 5 cm
\item ง) 8 cm
\end{enumerate}
\item \textbf{คลื่นเสียงที่มีความถี่สูงกว่า 20,000 Hz เรียกว่าอะไร}
\begin{enumerate}
\item ก) เสียงอินฟราซาวด์
\item ข) เสียงอัลตราซาวด์
\item ค) เสียงที่มนุษย์ได้ยิน
\item ง) โซนิกบูม
\end{enumerate}
\item \textbf{คลื่นวิ่งเป็นคลื่นที่มีรูปร่างและแอมพลิจูดคงที่ขณะเคลื่อนที่ ส่วนคลื่นนิ่งเกิดจากอะไร}
\begin{enumerate}
\item ก) การสะท้อนของคลื่นจากผนัง
\item ข) การแทรกสอดระหว่างคลื่นตกกระทบและคลื่นสะท้อนที่เคลื่อนที่ในทิศตรงข้าม
\item ค) การหักเหของคลื่น
\item ง) การเลี้ยวเบนของคลื่น
\end{enumerate}
\end{enumerate}
\end{document}